\section{Introduction}
% TODO: Motivate the problem, explain why existing approaches fall short,
% summarize the key idea and contributions, and preview the paper structure.

Modern datacenters host increasingly diverse distributed applications with heterogeneous network requirements. Even within a single application, different flows can have substantially different impacts on application performance. For example, in distributed machine learning training, parameter and tensor transfers may lie on different computation dependency paths: delaying some transfers can stall subsequent computation, while others can tolerate additional communication latency. Prior work, such as P3 and TicTac, has demonstrated that exploiting such dependencies to prioritize communication can substantially improve training performance. However, equal-share bandwidth allocation ignores differences in the importance of individual flows, strict priority scheduling can excessively delay lower-priority traffic, and static rate limits are poorly suited to changing traffic demands. **Fine-grained weighted bandwidth allocation provides a more flexible resource-sharing abstraction**, allowing applications to assign weights to individual flows and dynamically control their relative shares of network bandwidth to improve application performance.

Meanwhile, datacenter transports are evolving from traditional fixed-path forwarding toward increasingly flexible multipath transmission. This trend is particularly evident in large-scale AI training, where utilizing multiple network paths is critical to avoiding flow collisions, traffic hotspots, and communication slowdowns. For example, Meta's production experience highlights the limitations of conventional ECMP for AI training workloads. Google's Falcon employs congestion-aware packet-level multipathing, while Ultra Ethernet Transport (UET) explicitly supports packet spraying. OpenAI's Multipath Reliable Connection (MRC) uses adaptive packet spraying in large-scale AI supercomputing networks, and Alibaba Cloud's Stellar adopts RDMA packet spraying in large-scale AI clusters. Together, these developments illustrate a growing diversity of path-selection mechanisms, ranging from fixed-path forwarding and dynamic path switching to fine-grained packet spraying. **Bandwidth allocation should accommodate this diversity rather than remain tied to a particular path-selection mechanism.**

Unfortunately, providing consistent bandwidth allocation across different path-selection mechanisms is challenging. With conventional single-path forwarding, each flow traverses a relatively stable set of links, allowing its bandwidth competition to be analyzed in terms of the corresponding bottleneck resources. Under multipath transmission, however, packets belonging to the same logical flow may compete with different sets of flows on different paths or at different times. Packet spraying makes this challenge particularly pronounced, as consecutive packets can traverse different paths and experience different queueing delays and congestion conditions. Consequently, the bandwidth consumed by a logical flow cannot be readily controlled through an independent decision at a single path or bottleneck. Meanwhile, end hosts typically observe only aggregate end-to-end feedback, such as acknowledgments (ACKs) and round-trip times (RTTs), without directly knowing which paths individual packets traverse or how congestion is distributed across them. **The key challenge is therefore to control each flow's aggregate bandwidth share without relying on explicit path-level information.**

Existing bandwidth allocation mechanisms struggle to meet this requirement without incurring additional deployment or coordination costs. Switch-based schedulers, such as Weighted Fair Queueing (WFQ), can enforce weighted shares at individual output links, but are constrained by hardware queue and scheduling resources. Moreover, independent scheduling decisions at different links do not directly guarantee the aggregate bandwidth share of a logical flow spanning multiple paths. Centralized bandwidth allocators can coordinate allocations using network state, traffic demands, and routing information, but maintaining and processing such information becomes costly as traffic and path assignments change dynamically. Recent systems, such as Söze, have demonstrated lightweight decentralized bandwidth allocation, but still rely on switch-provided network telemetry and do not establish a unified allocation guarantee for packet-level multipathing. Other multipath-aware approaches often maintain separate congestion states or control information for individual paths or subflows, making them difficult to apply directly to packet spraying without explicit subflow separation. These limitations raise a fundamental question: **Can we provide fine-grained weighted bandwidth allocation across different transport modes using only standard end-to-end feedback, without path knowledge, additional switch support, or centralized coordination?**

We present **WAFT**, an end-host weighted bandwidth allocation system based on in-flight control. WAFT builds on a key insight: **bandwidth allocation need not operate on individual paths; instead, a single aggregate in-flight budget can regulate a logical flow's bandwidth share even when its packets traverse different network paths.** When competing flows experience comparable effective excess delays, their amounts of excess in-flight data are approximately proportional to their achieved bandwidth shares. WAFT exploits this relationship by mapping application-specified weights to per-flow excess in-flight budgets, enabling weighted bandwidth allocation without explicitly identifying packet paths or network bottlenecks.

Building on this insight, WAFT maintains a single aggregate in-flight window per flow and adjusts its window and pacing limit using only standard end-to-end feedback, such as ACKs and RTTs, while leaving path selection to existing load-balancing mechanisms. This design allows WAFT to employ the same control mechanism across single-path forwarding, dynamic path switching, and packet spraying, without maintaining separate allocation states for individual paths. In multipath environments, WAFT does not require identical path distributions across flows; instead, it leverages the capacity balancing and service consistency provided by the underlying load balancer to achieve path-oblivious bandwidth allocation. We further characterize how these service conditions affect allocation accuracy. **All allocation decisions are made independently at end hosts, without additional switch scheduling or telemetry support, centralized network-state collection, or global bandwidth computation**, substantially reducing deployment complexity and operational overhead.

We implement WAFT on **[implementation platform]** and evaluate it using **[testbed and simulation environments]**. Our evaluation covers **[validated single-path, dynamic path-switching, and packet-spraying scenarios]**, examining allocation accuracy and responsiveness under varying flow weights, contention levels, traffic loads, and load-balancing strategies. WAFT achieves **[allocation accuracy and convergence results]** under **[evaluated operating conditions]**, while incurring only **[end-host computation, memory, and control overhead]**. Furthermore, for **[application workloads]**, WAFT leverages application-specified flow weights to improve **[job or iteration completion time]** by **[measured improvement]**, while maintaining **[network utilization and latency results]**.

We make the following contributions:

- **A unified weighted bandwidth allocation mechanism.** We introduce a path-oblivious allocation approach that regulates aggregate bandwidth shares through flow-level excess in-flight budgets and characterize the network service conditions under which it operates.
- **A lightweight end-host implementation.** We design and implement WAFT, which performs decentralized weighted allocation using only standard ACK and RTT feedback, without network topology or packet-path information, additional switch assistance, or centralized coordination.
- **Evaluation across transport modes.** We evaluate WAFT's allocation accuracy, dynamic responsiveness, runtime overhead, and application-level benefits across single-path and multipath scenarios, and analyze the impact of underlying load-balancing behavior.